\section{Introduction}
\label{sec:introduction}

In the early 1990s the FBI held some 200 million fingerprint records on inked paper cards and
set out to digitize them. Scanned at 500 pixels per inch with 8 bits per pixel, a ten-print card
takes about 10 megabytes, so compression was a necessity. JPEG, the obvious candidate, codes the
image in $8\times 8$ blocks, and tests by the FBI and the UK Home Office found its tiling
artifacts unavoidable on fingerprints at compression ratios above about 10:1. Tom Hopper of the
FBI, with Jonathan Bradley and Christopher Brislawn of Los Alamos National Laboratory, designed
instead a codec built on a discrete wavelet transform, which has no blocks and preserves fine
detail such as sweat pores better at the same bit rate. The FBI adopted it in 1993 under the
name WSQ, for \emph{Wavelet Scalar Quantization} \cite{brislawn1995,bradley1994}.

WSQ remains the standard compression for 500 ppi fingerprint images in ANSI/NIST-ITL
transactions (record types 4 and 14) and in the FBI's EBTS. A typical 500 ppi image compressed at
0.75 bits per pixel shrinks by a factor of about 15.

\subsection{What the specification fixes}

The specification \cite{wsq31} has three parts. Part~1 defines the decoder, the compressed
data format and a \emph{class} of encoders. Part~2 defines how implementations are tested for
compliance. Part~3 fixes the parameters of the one encoder the FBI approves, called
\emph{encoder number two} since version~3.1 of 2010.

The split matters. Everything the decoder needs travels inside the file: the wavelet filters,
the quantizer bin widths and the Huffman tables. A decoder is therefore generic: it must accept
any filter bank of the allowed kinds, any bin widths and any Huffman tables. An encoder, on the
other hand, only has to produce valid streams with acceptable accuracy, and an FBI-compliant
encoder must in addition choose its parameters exactly as Part~3 prescribes. Version~3.1 changed
only that choice (the way subband variances are estimated, Section~\ref{sec:variance}); the
decoder is the same as in 1993.

\subsection{Overview}

Figure~\ref{fig:pipeline} shows the pipeline. The encoder
\begin{enumerate}
  \item shifts and scales the image to a signed range (Section~\ref{sec:normalization});
  \item decomposes it with a two-dimensional symmetric wavelet transform into 64 subbands, each
        holding one band of horizontal and vertical spatial frequencies
        (Section~\ref{sec:transform});
  \item quantizes each subband with a uniform scalar quantizer whose step depends on the
        subband's variance and on the target bit rate (Section~\ref{sec:quantization});
  \item codes the quantizer indices with Huffman codes designed for the image
        (Section~\ref{sec:entropy});
  \item writes everything, tables included, as a sequence of marker segments in the style of
        JPEG (Section~\ref{sec:format}).
\end{enumerate}
The decoder undoes each step in reverse order (Section~\ref{sec:decoding}). Only the
quantization loses information; the transform is invertible and the entropy coding is
lossless. Section~\ref{sec:compliance} describes how implementations are tested and reports
how the procedures in this document compare with the NIST reference implementation.
Section~\ref{sec:implementation} collects implementation notes.

\begin{figure}[ht]
  \centering
  \begin{tikzpicture}[node distance=5mm and 4.5mm,
    stage/.style={box,minimum width=18mm,font=\small,align=center},
    side/.style={font=\footnotesize,align=center,text=black!75}]
  % encoder
  \node[font=\small] (img) {$I(m,n)$};
  \node[stage,right=of img] (norm) {normalize};
  \node[stage,right=of norm] (swt) {wavelet\\transform};
  \node[stage,right=of swt] (quant) {quantize};
  \node[stage,right=of quant] (huff) {Huffman\\encode};
  \node[font=\small,right=of huff] (out) {stream};
  \draw[->] (img) -- (norm);
  \draw[->] (norm) -- (swt);
  \draw[->] (swt) -- node[above,font=\footnotesize] {$a_k$} (quant);
  \draw[->] (quant) -- node[above,font=\footnotesize] {$p_k$} (huff);
  \draw[->] (huff) -- (out);
  % side information
  \node[side,below=4mm of norm] (sM) {$M,\ R$\\\marker{SOF}};
  \node[side,below=4mm of swt] (sF) {$h_0,\ h_1$\\\marker{DTT}};
  \node[side,below=4mm of quant] (sQ) {$Q_k,\ Z_k,\ C$\\\marker{DQT}};
  \node[side,below=4mm of huff] (sH) {tables\\\marker{DHT}};
  \foreach \s/\t in {sM/norm,sF/swt,sQ/quant,sH/huff} \draw[densely dotted] (\t) -- (\s);
  % decoder
  \node[stage,below=4mm of sM] (dnorm) {denormalize\\and round};
  \node[stage,below=4mm of sF] (dswt) {inverse\\transform};
  \node[stage,below=4mm of sQ] (dquant) {dequantize};
  \node[stage,below=4mm of sH] (dhuff) {Huffman\\decode};
  \node[font=\small] at (out |- dhuff) (in) {stream};
  \node[font=\small] at (img |- dnorm) (rec) {$\hat I(m,n)$};
  \foreach \s/\t in {sM/dnorm,sF/dswt,sQ/dquant,sH/dhuff} \draw[densely dotted] (\s) -- (\t);
  \draw[->] (in) -- (dhuff);
  \draw[->] (dhuff) -- node[above,font=\footnotesize] {$p_k$} (dquant);
  \draw[->] (dquant) -- node[above,font=\footnotesize] {$\hat a_k$} (dswt);
  \draw[->] (dswt) -- (dnorm);
  \draw[->] (dnorm) -- (rec);
\end{tikzpicture}

  \caption{The WSQ encoder (top) and decoder (bottom). The labels between the rows are the
  parameters the encoder writes into the stream and the marker segments that carry them.}
  \label{fig:pipeline}
\end{figure}

\subsection{Notation}

An image has $Y$ lines (rows) and $X$ samples per line (columns), the names the
specification uses in the frame header. A pixel is $I(m,n)$ with row $m$ and column $n$,
both counted from zero at the upper left corner. The 64 subbands are numbered
$k = 0,\dots,63$; subband $k$ has $X_k$ columns and $Y_k$ rows and holds the wavelet
coefficients $a_k(m,n)$. Quantization maps them to integer indices $p_k(m,n)$ with bin width
$Q_k$ and zero-bin width $Z_k$. The analysis filters are $h_0$ (lowpass) and $h_1$
(highpass); the synthesis filters are $f_0$ and $f_1$. For a real $x$, $\floor{x}$ is the
largest integer not above $x$ and $\ceil{x}$ the smallest integer not below it. Hexadecimal
byte values are written \hex{FF}.
